%%%%%%%%%%%%%%%%%%%%%%%%%%%%%%%%%%%%%%%%%%%%%%%%%%%
%
%  New template code for TAMU Theses and Dissertations starting Spring 2021.  
%
%
%  Author: Thesis Office
%  
%  Last Updated: 1/13/2021
%
%%%%%%%%%%%%%%%%%%%%%%%%%%%%%%%%%%%%%%%%%%%%%%%%%%%
%%%%%%%%%%%%%%%%%%%%%%%%%%%%%%%%%%%%%%%%%%%%%%%%%%%%%%%%%%%%%%%%%%%%%
%%                           ABSTRACT 
%%%%%%%%%%%%%%%%%%%%%%%%%%%%%%%%%%%%%%%%%%%%%%%%%%%%%%%%%%%%%%%%%%%%%

\chapter*{ABSTRACT}
\addcontentsline{toc}{chapter}{ABSTRACT} % Needs to be set to part, so the TOC doesn't add 'CHAPTER ' prefix in the TOC.

\pagestyle{plain} % No headers, just page numbers
\pagenumbering{roman} % Roman numerals
\setcounter{page}{2}

% Modularity is an important aspect of solid modeling, with modularity the design space can be enumerated and the manufacturing costs can be reduced. When we try to apply modularity to woven structures, the weave threads admit a symmetry that decomposes the space into modules. This partition creates a piecewise hyperbolic surface where the surface parameters can be parametrized by parametric curves.

% The 

% Modularity and symmetry are fundamental properties of woven structures that govern their geometric versatility and mechanical behavior. While traditional studies have focused on planar weaves, recent advances demonstrate the untapped potential of volumetric weave structures that extend fabric-like topologies into three-dimensional space. This thesis bridges the gap between planar and volumetric weave design by unifying geometric modularity, algebraic construction, and lattice symmetry into a cohesive framework.

% First, we explore modular space-filling threads, demonstrating that a single yin-yang type congruent tile partitioned from a cube and foldable from planar materials can generate infinite families of two-way, two-fold woven structures through translation and rotation. This modular approach eliminates the need for secondary fillers, enabling self-supporting volumetric fabrics. Next, we introduce an algebraic framework for nonwoven volumetric textiles, which generalizes planar nonwoven structures into 3D. These structures mimic woven behavior while exhibiting unprecedented material properties, including representations of infinite regular polyhedra and flexible frameworks. Finally, we formalize volumetric weaves through the isometries of Bravais lattices, defining a systematic design space for cP that surpasses the complexity of planar weaves by an order of magnitude.

% Together, these contributions establish a foundation for designing symmetry-driven and modular volumetric weaves, with implications for engineered materials, composite structures, and digital fabrication. By leveraging geometric invariance, algebraic parametrization, and lattice isometries, this work expands the taxonomy of weaves beyond planar constraints, unlocking new paradigms for functional material design.

% Weaving is one of the oldest human activities, it is a process of creating arrangements of filaments via interlacing different threads. Weaving emphasizes one of the greatest knowledge of the humankind, the whole is a greater than the sum of parts, woven objects shows various material properties and functional use such that the same material propertoes are not found in the individual threads. The use of weaves in daily life varies from baskets to fabrics, the higher technology knowledge required by fields like construction and high performance composites also benefits from weaves as well. Weaves are also a part of the nature, the similar material properties have been observed between weaves and the elastic biological tissues.  

Weaving is one of humanity’s oldest activities, involving the creation of structures by interlacing threads or filaments. It exemplifies a fundamental principle: the whole is greater than the sum of its parts. Woven objects display material properties and functions not present in the individual threads. Applications of weaving range from everyday objects such as baskets and fabrics to advanced fields such as construction and high-performance composites. Moreover, similar structural principles appear in nature, where woven-like arrangements are observed in elastic biological tissues.

% Woven structures are arrangements of interlaced threads or filaments that form ordered, repeating patterns. Beyond textiles and clothing, woven structures have long attracted attention for their mathematical regularity and structural efficiency. They appear in contexts ranging from architectural design to engineered composites, and more recently in digital fabrication and materials science. Their appeal lies not only in aesthetic elegance but also in their ability to achieve lightweight, flexible, and mechanically resilient architectures. Strength of traditional fabrics is anisotropic, manifesting itself primarily in the direction of the thread orientations. That is if the threads are arranged in a plane. To extend the use and value of textiles into industrial and engineering applications, which typically require strength in more than two directions, textile engineers were challenged to develop stronger thread architectures. So they come up with different ways to manufacture three-dimensional woven textiles. From the manufacturing point of view the easiest way to realize such structures were to manufacture the textile layer by layer and connect different layers by straight threads. Even though the manufacturing of these textiles were developed, there is a gap in designing and understanding the three dimensional weaves. In this thesis we address this gap of design space. Like weaves, this thesis starts with a single thread and weaves a connection between the different concepts to finally develop a framework to design three dimensional woven textiles.

Woven structures consist of interlaced threads forming ordered, repeating patterns. Beyond textiles, they are valued for their structural efficiency and appear in architecture, engineered composites, and digital fabrication. Traditional fabrics are anisotropic, exhibiting strength primarily along thread directions in a plane. Extending textiles to industrial and engineering applications—where strength is needed in multiple directions—led to the development of three-dimensional woven textiles. These volumetric weaves, which generalize fabrics into three dimensions, offer new opportunities for material design but currently lack a unified theoretical framework. From a manufacturing perspective, the simplest approach has been to build these textiles layer by layer, connecting layers with straight threads. Despite advances in manufacturing, there remains a gap in systematically designing and understanding three-dimensional woven structures. This thesis addresses that gap, starting from a single thread and connecting concepts to develop a framework for designing volumetric woven textiles.

% Traditional studies of weaving focus primarily on planar arrangements: two-dimensional patterns repeated across a surface. These weaves can be represented by relatively simple algebraic rules and have been extensively catalogued in both mathematics and textile engineering. Yet, limiting weaves to planar surfaces overlooks the rich possibilities of three-dimensional interlacement, where threads can extend through volume rather than across a plane. Volumetric weaves, woven analogues of fabrics extended into 3D, open new opportunities for material design, but lack a unified theoretical foundation.

% Traditional studies focus on planar weaving: two-dimensional patterns repeated across a surface. These weaves can be described using relatively simple algebraic rules and have been extensively catalogued in mathematics and textile engineering. However, restricting weaving to planar surfaces overlooks the possibilities of three-dimensional interlacement, where threads extend through volume rather than a plane. Volumetric weaves—three-dimensional analogues of fabrics—offer new opportunities for material design but currently lack a unified theoretical foundation.

% To understand the the weaves in three dimension, we first start the problem from an earlier stage. In mathematics, weaves in 2-dimensions, i.e. planar weaves are encoded with either a square grid or a triangular grid depending on the type of weave. So to understand the whole, the system has been modularized to tiling of the 2D space. In the first part of this thesis we approach to the problem of space-filling threads. While previous work produced space-filling threads, these were fixed and non-reconfigurable.  By contrast, the modular tiles allows easy reconfiguration through translation and rotation, enabling the generation of infinite families of bidirectional, planar weaves. Beyond reducing the need for filler elements, this modularity highlights the role of symmetry in fabric design, motivating the extension of these principles into 3D weaves.

To study three-dimensional weaving, we begin by revisiting planar weaves. In mathematics, planar weaves can be represented using square or triangular grids depending on the pattern. To better understand these structures, we modularize the system by tiling the 2D space. The first part of this thesis addresses the problem of space-filling threads. Previous approaches produced fixed, non-reconfigurable threads, whereas our modular tiles allow reconfiguration through translation and rotation, generating infinite families of bidirectional planar weaves. Beyond reducing the need for filler elements, this modularity highlights the role of symmetry in fabric design, providing a foundation for extending these principles into volumetric weaving.

Second, we develop an algebraic framework for nonwoven volumetric textiles. Nonwoven structures differ from woven ones in that they do not strictly follow the over–under interlacing rule; instead, they are networks of paths that still achieve fabric-like connectivity. By extending these ideas into three dimensions, we show how we can generate symbols such that we can identify families of structures by rules rather than by direct construction, a process we call algebraic parametrization. This approach allows the creation of configurations that mimic the mechanical behavior of woven fabrics while also realizing structures not possible in planar weaves.

Finally, we formalize volumetric weaves using the isometries of Bravais lattices. Extending the modularity concept to three dimensions allows us to define volumetric weaves as assemblies of repeatable modules, analogous to the regularity observed in crystals. This lattice-based framework provides a systematic design space for volumetric weaves, enabling locally configurable structures that go beyond current understanding and achieve complexity far exceeding that of planar weaves.

Together, these studies establish a framework for symmetry-driven, modular design of volumetric weaves. By integrating geometric invariance (properties unchanged under symmetry transformations), algebraic parametrization (symbolic generation of structure families), and lattice isometries (crystalline symmetries as organizing principles), this work provides a unified foundation for three-dimensional weaving. The framework has applications in engineered materials, composite structures, and digitally fabricated architectures. More broadly, it demonstrates how weaving, traditionally a planar craft, can be extended into a systematic science of volumetric interlacement, expanding possibilities for both functional and aesthetic material design.

% \indent This is the first numbered page, lower case Roman numeral (ii). Page numbers are outside the prescribed margins, at the bottom of the page and centered; everything else is inside the margins.No bold on this page except the heading ABSTRACT if all major headings are bold. \emph{This \LaTeX ~ template applies to this exception}).

% Text begins two double spaces below the major heading. The Abstract should be no more than 350 words. Vertical spacing is double spaced. (\emph{This \LaTeX ~ template applies double space for this ABSTRACT.}) The margin settings and text alignment should be consistent throughout the document. There should be no numbered references or formal citations in ABSTRACT.

% The content of this ABSTRACT provides a complete, short snapshot of the research, addressing the purpose, methods, results, and conclusions of the document. As a result, it should stand alone without any formal citations or references to chapters/sections of the work. To accommodate with a variety of online database, images, complex equations, or Greek letters/symbols should also be avoided.This should be no longer than 350 words.

% The next pages are Dedication, Acknowledgments, Contributors and Funding Sources, and Nomenclature. Contributors and Funding Sources is required, the rest are optional.


 

\pagebreak{}
